\section{Introduction}
\label{sec:benchmarks_motivation}

\Cref{ch:novelty_open_world} characterized open-world environments as settings in which an agent's internal models are incomplete with respect to the environment; \cref{ch:foundations} formalized novelty as a mismatch between those models and observed outcomes. Principled evaluation in this setting must characterize how agents detect, localize, and recover from mismatch.

Evaluating agents under mismatch presents challenges that do not arise in closed-world settings. In standard benchmarks, the environment is fixed and fully specified, and performance can be measured with respect to a known objective. In contrast, open-world environments are characterized by the possibility of unanticipated changes: objects may appear or disappear, action effects may change, and task structure may evolve. These changes induce mismatch between the agent's expectations and the environment, but the space of possible mismatches is unbounded. As a result, evaluation cannot rely on a fixed set of scenarios, but must instead provide mechanisms for generating controlled and interpretable variations in the environment.

The crafting scenario introduced in Chapter~\ref{ch:foundations} illustrates this challenge. When the agent encounters the axe novelty---where trees become unbreakable without a new tool---its existing plan fails due to a mismatch between the expected and actual effects of the \texttt{break} operator. Evaluating the agent's behavior in this setting requires more than measuring task success: it requires quantifying how quickly the agent detects the mismatch, how severely its performance degrades, and how efficiently it recovers. More generally, evaluation must capture the dynamics of adaptation, not just its final outcome.

Open-world research also requires testbeds that support agents with different representational and decision-making commitments. Comparisons among symbolic planning, reinforcement learning, and hybrid architectures can reveal how an agent's internal structure affects novelty detection, performance degradation, and recovery. These comparisons require metrics for the course of adaptation in addition to final task success. The benchmarks introduced in this chapter provide these capabilities while leaving the choice of agent architecture and adaptation method open.

To evaluate open-world agents along these dimensions, we require environments that can systematically introduce mismatch, expose its effects at both the planning and execution levels, and support different agent paradigms. This chapter presents two such environments. \NovelGridworlds{} introduces a structured taxonomy of novelties and evaluation metrics that capture performance degradation and detection. \NovelGym{} extends this framework with a formal environment representation, composable transformations that induce mismatch, and metrics that explicitly characterize adaptation dynamics. Together, they form a general-purpose suite for developing, studying, and evaluating open-world agents, including sophisticated hybrid planning--learning architectures. The experiments in this chapter demonstrate this capability across several agent classes, and \cref{ch:rapid_learn} later uses \NovelGridworlds{} as an evaluation setting for \RAPidLearn{}. The environments, protocols, and metrics are intended as research infrastructure for the broader community and are independent of any single adaptation method.
